\chapter{Vector Representation in SVG Generation}\label{ch:representation}

SVG generation models differ not only in their architectures and 
training objectives but also in how they represent vector graphics 
internally during the generation process. While all approaches 
ultimately produce SVG files as output, the internal representation 
used during learning or optimization varies considerably across 
methods. This variation has direct consequences for how geometry is 
parameterized, how supervision signals are applied, and what 
structural properties the resulting graphics have.

In the context of this thesis, representation refers to the 
intermediate form in which a model encodes vector geometry, not the 
final output file but the internal structure through which shapes are 
created, manipulated, and decoded. A vector graphic can be represented 
in several fundamentally different ways. It can be described as a 
collection of explicit geometric primitives with directly optimizable 
parameters. It can be encoded as an ordered sequence of drawing 
commands that procedurally construct the image. It can be defined 
implicitly through the weights of a neural network that maps indices 
or coordinates to shape properties. Or it can be embedded into a 
compact latent space from which the model samples and decodes 
structured output. Each of these formulations encodes the same 
underlying geometry differently, and this difference shapes how the 
model learns, how gradients propagate, and what the generated SVG 
looks like structurally.

Representation is therefore not a neutral implementation detail. It 
determines whether a model can produce clean hierarchical structure, 
how easily the output can be edited after generation, and how well the 
method scales to complex visual content. These properties are central 
to the comparative analysis in this thesis.

This chapter examines four primary representation families identified 
across the selected models: parametric, sequential, implicit neural, 
and latent. For each family, the analysis describes how it works 
technically, where it originates, and what structural tendencies it 
introduces in SVG generation.



\section{Parametric Representations}

Parametric representations describe vector graphics as collections of 
explicit geometric primitives, specifically paths defined by Bézier 
control points, stroke attributes, and fill colors. These parameters 
can be optimized directly using gradient-based methods or predicted by 
a neural network. The representation closely matches the SVG format 
itself, making the internal state of the model directly interpretable 
in geometric terms.

The use of parametric representations in neural SVG generation became 
practical with the introduction of DiffVG~\cite{liDifferentiableVectorGraphics2020a}, 
which made the rasterization of Bézier primitives differentiable. 
Before this, gradients could not flow from pixel-space losses back to 
vector parameters. DiffVG enabled models to optimize control points 
and colors directly against image-based or semantic objectives, which 
is the foundation on which LIVE~\cite{maLayerwiseImageVectorization2022}, 
SAMVG~\cite{zhuSAMVGMultistageImage2023}, 
VectorFusion~\cite{jainVectorFusionTexttoSVGAbstracting2022}, and 
SVGDreamer~\cite{xingSVGDreamerTextGuided2024} are all built.

In practice, these models initialize a fixed set of vector primitives 
and iteratively refine their parameters so that the rendered image 
aligns with a target signal. The initialization strategy varies: LIVE 
uses a layer-by-layer approach that progressively adds paths, while 
SVGDreamer's SIVE module uses diffusion model attention maps to 
determine where primitives should be placed before optimization 
begins~\cite{xingSVGDreamerTextGuided2024}. Despite these differences in 
initialization, all of these methods share the same underlying 
representational commitment, namely that shapes are stored as explicit 
geometric parameters that are modified directly during optimization.

The main limitation of parametric representations is that complexity 
is handled by increasing the number of primitives. When a model needs 
to capture fine visual detail, it tends to introduce many overlapping 
or redundant paths rather than increasing the expressiveness of 
individual ones. LIVE, for example, uses 160 paths across five layers 
to represent a single image~\cite{maLayerwiseImageVectorization2022}. At 
this scale, the output is theoretically accessible for editing since 
each path is a geometric object, but in practice the number and 
entanglement of paths makes meaningful manual editing very difficult. 
SVGDreamer++ explicitly identifies path entanglement as a failure mode 
of optimization-based parametric methods~\cite{xingSVGDreamerTextGuided2024}.

\section{Sequential Representations}

Sequential representations treat vector graphics as ordered sequences 
of drawing commands. Rather than storing shapes as independent 
geometric objects, the model generates a series of commands that 
describe how the graphic is constructed step by step. This mirrors the 
structure of SVG path data directly.

The use of sequences to represent drawings originates in work on 
sketch generation. Ha and Eck's SketchRNN~\cite{haNeuralRepresentationSketch2017} 
demonstrated that recurrent neural networks could learn to generate 
coherent hand-drawn sketches by modeling pen stroke sequences 
autoregressively. This established the idea that a drawing could be 
treated as a variable-length sequence and modeled with the same 
architectures used for language. DeepSVG~\cite{carlierDeepSVGHierarchicalGenerative2020} 
extended this approach to structured SVG icons by combining a 
hierarchical VAE with a transformer decoder that generates sequences 
of SVG path commands. DeepIcon~\cite{bingDeepIconHierarchicalNetwork2024} follows the 
same sequential formulation, learning to generate icon SVGs from 
command-level supervision.

Sequential representations allow the use of powerful sequence modeling 
architectures and can capture drawing patterns present in large 
datasets. However, they have known scaling limitations. Long command 
sequences are difficult to model reliably, and prediction errors early 
in the sequence propagate through the entire graphic. More 
fundamentally, a procedural command sequence does not encode explicit 
spatial hierarchy, meaning the model generates commands in order but 
nothing in the representation directly encodes which commands belong 
to which semantic region. Editability is therefore limited even when 
visual reconstruction quality is high.

\section{Implicit Neural Representations}

Implicit neural representations encode vector graphics through the 
weights of a neural network rather than through explicit geometric 
parameters. Instead of storing control points or commands, the model 
learns a function that maps indices or coordinates to shape properties 
such as geometry, color, and layer assignment. Shapes are defined 
implicitly by evaluating this function.

This approach is directly inspired by Neural Radiance Fields 
(NeRF)~\cite{mildenhallNeRFRepresentingScenes2020}, which demonstrated that a compact 
MLP could represent a full 3D scene implicitly by mapping spatial 
coordinates to color and density values. NeuralSVG explicitly 
acknowledges this connection: each SVG is represented as a set of 
indices where each index corresponds to a single shape, and a small 
MLP learns to map these indices to Bézier control points and fill 
colors~\cite{polaczekNeuralSVGImplicitRepresentation2025}. The network weights are the 
representation, not the shapes themselves. This also allows NeuralSVG 
to adopt an ordered representation strategy, where shapes are assigned 
to fixed indices during training to ensure stable correspondence 
between the generated output and the supervision signal.

NIVeL follows a related paradigm, representing layered vector graphics 
through neural functions optimized against image-based supervision. 
Rather than predicting individual path parameters as separate outputs, 
NIVeL encodes the entire layered structure implicitly, where each layer 
is defined by a continuous neural function over the image 
domain~\cite{thamizharasanNIVeLNeuralImplicit2024}. This enables the model to produce 
clean layered decompositions without requiring explicit path 
initialization.

Implicit representations offer strong expressive power and can produce 
outputs with coherent global structure because the entire graphic is 
encoded in a single set of network weights. The main limitation is 
interpretability. Individual shapes are not stored as separate 
accessible objects, and converting the implicit output into clean, 
directly editable SVG primitives requires an additional extraction 
step. The quality of this conversion determines how usable the output 
is in practice.

\section{Latent Representations}

Latent representations introduce an abstraction layer by encoding 
vector graphics as compact embeddings in a learned latent space. 
Rather than representing shapes directly, the model learns a space in 
which structured SVG content can be sampled, interpolated, or 
denoised. Generation then occurs by decoding these embeddings into 
explicit vector primitives or command sequences.

The foundational mechanism is the variational autoencoder 
(VAE)~\cite{kingmaAutoEncodingVariationalBayes2022}, which trains an encoder to compress 
structured data into a continuous latent space and a decoder to 
reconstruct it. DeepSVG~\cite{carlierDeepSVGHierarchicalGenerative2020} applies this to 
SVG generation by encoding path command sequences into latent vectors 
and training a transformer decoder to reconstruct them. The latent 
space captures structural regularities across the training dataset, 
enabling the model to sample new SVGs by decoding points sampled from 
this space.

Latent diffusion extends this idea by applying diffusion-based 
generative modeling directly in the latent space of a pretrained 
vector encoder~\cite{rombachHighResolutionImageSynthesis2021}. SVGFusion operates 
in this way, treating the latent vector representation as the space in 
which the diffusion process runs. T2V with Neural Path Representation 
takes a different approach, arguing that global SVG-level latent spaces 
are too constrained to capture diverse path combinations, and instead 
learns a neural path representation at the individual path 
level~\cite{zhangTexttoVectorGenerationNeural2024}.

Latent representations improve scalability and generative diversity 
compared to direct parametric or sequential approaches. The trade-off 
is that edits in latent space affect the decoded output in ways that 
are not always predictable at the level of individual shapes. 
Fine-grained geometric control is harder to achieve than with 
parametric representations because the relationship between latent 
dimensions and specific visual properties is not explicit.

\section{Hybrid Representations}

In practice, many SVG generation models combine elements from more 
than one representation family. These combinations reflect the fact 
that no single representation fully addresses the competing 
requirements of structural correctness, editability, scalability, and 
visual realism.

Optimization-based diffusion methods such as VectorFusion and 
SVGDreamer use parametric primitives as the direct optimization target 
while relying on a latent diffusion model as the semantic prior. The 
parametric layer provides geometric interpretability, while the latent 
model provides semantic coherence that pure geometric optimization 
cannot achieve alone. SVGFusion takes a similar combined approach, 
encoding path-level structure into latent embeddings before applying 
diffusion, effectively combining latent and sequential elements in a 
single pipeline.

LayerTracer~\cite{songLayerTracerCognitiveAlignedLayered2025} represents a different kind of 
hybrid strategy. Rather than combining representation families at 
inference time, it integrates them during training. The model is 
trained on the Serpentine dataset, which contains step-by-step human 
SVG construction sequences, meaning the training data encodes a 
sequential procedural representation. However, the output the model 
learns to produce is a layered parametric SVG, where each layer 
corresponds to a meaningful visual component. The result is that 
editability becomes a learned property rather than an emergent one: 
because the model has been trained on human drawing sequences, it 
learns to produce outputs that are structured the way a human designer 
would structure them, with clean layer separation and manageable path 
counts.

NeuralSVG also combines ideas across families. Its core representation 
is implicit, with an MLP mapping shape indices to geometric parameters, 
but it incorporates an ordered assignment mechanism originally 
introduced in DeepSVG to ensure that shape indices correspond 
consistently to visual regions acrosstraining~\cite{polaczekNeuralSVGImplicitRepresentation2025}. This ordered structure imposes 
a form of explicit hierarchy on what is otherwise a purely neural 
representation.

These hybrid strategies highlight that representation design in SVG 
generation is less about choosing one family and more about deciding 
which trade-offs to prioritize. The following chapters examine how 
these trade-offs manifest in the supervision strategies and qualitative 
outputs of the selected models.


\section{Representation and Editability}

Editability is one of the defining advantages of vector graphics and 
a central requirement in real design workflows. Unlike raster images, 
SVG graphics are expected to support direct manipulation of individual 
elements, including moving shapes, adjusting colors, reorganizing 
layers, or modifying geometry without affecting unrelated parts of the 
graphic. Representation plays a central role in determining how 
editable a generated SVG is.

Parametric representations offer the most direct path to editability 
because each element corresponds to explicit geometric parameters. 
However, this advantage depends heavily on path count and geometric 
organization. When models such as LIVE or SAMVG introduce large 
numbers of paths to match visual detail, the output becomes 
theoretically editable but practically difficult to work with. 
Sequential representations such as those used in DeepSVG and DeepIcon 
produce procedurally structured output, but editing a single visual 
element requires locating its corresponding command segment, which is 
not always straightforward. Implicit representations shift geometry 
into neural parameter space, meaning individual shapes are not 
directly accessible until the network output is converted into 
discrete primitives. How clean this conversion is determines whether 
the final SVG supports practical editing.

An important observation is that editability is not determined by 
representation type alone. Initialization strategy, supervision 
signals, and structural priors all contribute. Methods that 
incorporate segmentation, attention-based masking, or layer-aware 
initialization tend to produce outputs with clearer object boundaries 
regardless of their underlying representation. LayerTracer is a clear 
example: by training on human construction sequences, it learns to 
produce layer-separated outputs as a direct consequence of its 
training data rather than through explicit structural constraints.

This suggests that editability emerges from the interaction between 
representation, optimization strategy, and structural design rather 
than from representation alone. Representations that maintain explicit 
primitives support local editing more naturally, while representations 
that emphasize abstraction improve generative flexibility at the cost 
of direct control. Hybrid approaches attempt to balance this by 
preserving layered structure while allowing flexible optimization or 
semantic guidance.


\section{Structural Properties Emerging from Representation}

Beyond editability, representation choices directly influence the 
structural properties of generated SVG graphics. Structural quality 
refers to how well a vector graphic preserves meaningful organization, 
including hierarchy, grouping, layering, and geometric consistency. 
While visual similarity can be achieved through many small primitives, 
structural properties determine whether the SVG remains interpretable 
and aligned with design practices.

Parametric optimization-based approaches frequently illustrate the 
tension between visual accuracy and structural quality. Because these 
methods refine primitives to minimize image-space loss, they tend to 
introduce redundant paths or fragmented geometry to capture fine 
visual detail. This improves reconstruction quality but can result in 
overlapping shapes and inconsistent grouping. Sequential 
representations produce structure implicitly through command order 
rather than explicit hierarchy, which can leave structural 
relationships ambiguous even when command prediction is accurate. 
Implicit neural representations allow geometry to emerge from 
continuous neural functions, enabling smoother global organization and 
layered outputs, but structural clarity depends on how carefully the 
conversion into discrete SVG elements is handled. Latent 
representations capture structural regularities from training data, 
which can lead to more coherent outputs compared to purely 
optimization-based methods, but clean hierarchy is not guaranteed 
unless structural signals are explicitly modeled during training.

A key observation is that structural properties in most models emerge 
indirectly. The majority of approaches optimize visual objectives, 
with structure appearing as a byproduct rather than a primary 
optimization target. This explains why many methods achieve strong 
visual fidelity but still produce SVGs with redundant primitives or 
weak hierarchy. Initialization strategies partially address this: 
methods that begin from segmentation masks, attention priors, or 
layer-aware starting points tend to produce cleaner organization 
because they reduce the need for excessive path refinement. 
Representation defines the space in which structural organization can 
appear, but initialization and supervision determine whether that 
structure remains coherent in the final output.


\section{Summary}

This chapter examined representation as a central analytical dimension 
for understanding SVG generation methods. Although all models 
ultimately produce SVG files, the way vector graphics are encoded 
during generation introduces fundamental differences in structure, 
editability, and visual behavior. Representation therefore acts as a 
unifying perspective for interpreting the design decisions underlying 
different approaches.

Four primary representation families were identified: parametric, 
sequential, implicit neural, and latent. Parametric representations 
provide direct geometric control and strong interpretability but may 
produce fragmented structures under complex optimization. Sequential 
representations capture procedural drawing behavior but struggle with 
scalability and explicit hierarchy. Implicit neural representations 
enable flexible topology and expressive layered outputs while shifting 
optimization into neural parameter space. Latent representations offer 
scalable generative modeling and global style control at the cost of 
fine-grained geometric control.

Across these families, representations that maintain explicit 
primitives tend to support local editing and structural transparency, 
whereas representations that emphasize abstraction improve generative 
flexibility. Hybrid approaches combine explicit structure with learned 
or semantic guidance to navigate this trade-off. Editability and 
structural quality rarely emerge from representation alone but depend 
on the interaction between representation, initialization strategy, 
and supervision signals.

These observations provide the foundation for the next analytical 
dimension of the thesis. While representation defines how vector 
graphics are encoded, supervision determines how models learn to shape 
that representation. The following chapter examines supervision 
signals and guidance strategies, analyzing how they interact with 
representation choices to influence the quality, structure, and 
behavior of generated SVG graphics.